\section{Root of Unity Filters}

\begin{theorem}
For any polynomial or convergent series \(F(x) = a_0 + a_1x + a_2x^2 + a_3x^3 + \dots\), 
you can filter out the sum of every $n$-th coefficient $(a_0 + a_n + a_{2n} + \dots)$ 
using the $n$-th roots of unity:
\[\sum _{k=0}^{\infty }a_{kn}=\frac{1}{n}\sum _{j=0}^{n-1}F(\omega ^{j})\]
where \(\omega = e^{\frac{2\pi i}{n}}\) is a primitive n-th root of unity.
\end{theorem}

To see exactly why this cancellation happens, we can analyze the sum using the geometric series formula.
Let the sum be \(S = \sum_{k=0}^{n-1} (\omega^m)^k\), where \(\omega = e^{\frac{2\pi i}{n}}\).\\

Case 1: When \(m\) is a multiple of \(n\)\\
If \(m\) is a multiple of \(n\) (so \(m = c \cdot n\) for some integer \(c\)):

Each term becomes \((\omega^n)^c\).
Since \(\omega^n = 1\), each term equals \(1^c = 1\).
The sum adds \(1\) exactly \(n\) times.
Result: \(S = 1 + 1 + \dots + 1 = n\).\\

Case 2: When \(m\) is not a multiple of \(n\)\\
When \(m\) is not a multiple of \(n\), the common ratio \(r = \omega^m\) is not equal to \(1\).
We apply the finite geometric series formula, \(S = \frac{1 - r^n}{1 - r}\): [1]\\

Substitute the ratio: \(S = \frac{1 - (\omega^m)^n}{1 - \omega^m}\)
Swap the exponents: \(S = \frac{1 - (\omega^n)^m}{1 - \omega^m}\)
Evaluate the numerator: Since \(\omega^n = 1\), the numerator becomes \(1 - 1^m = 0\).
Ensure non-zero denominator: Because \(m\) is not a multiple of \(n\), \(\omega^m \neq 1\), so the denominator is not zero.
Result: \(S = \frac{0}{\text{non-zero}} = 0\).\\

Visualizing the Cancellation
Geometrically, when \(m\) is not a multiple of \(n\), the terms \(\omega ^{km}\) represent vectors pointing symmetrically in all directions around the unit circle. Because of this perfect rotational symmetry, their vectors balance each other out completely, pulling the net sum directly to the origin \((0,0)\). [1]\\


%Theorem: (Root of Unity Filter)
%Define \(?=e^{2�i/n}\)
 %for a positive integer \(n\). For any polynomial \(F(x)=a+0+a_1x+a_2x^2+�(where we take a_k=0 if  k>d eg(F)), the sum a0+an+a2n+... is given by a0+an+a2n+?=1 n (F (1)+F (?)+?+F(?n?1))
%Proof: Let sk=1+?k+?+?(n?1)kIf n divides k, then ?k=1 and so sk=1+1+1?+1=n otherwise sk=1??nk1??k=0. So F(1)+F(?)+?+F(?n?1)= a0 s0+a1s1+a2s2+?=n(a0+an+a2n+�)
%Divide the both sides of the equation by n and the proof is complete.
%Source of my knowledge: http://zacharyabel.com/papers/Multi-GF_A06_MathRefl.pdf
%There are some examples also which may help you. Please have a look at Problem 2 on page 3.

The roots of unity filter (or roots of unity trick) is a mathematical technique used in generating functions and combinatorics to isolate and sum specific polynomial coefficients, such as every $n$-th term. When applied to a product of polynomials, it extracts those terms whose exponents are periodic.\\

The core of the filter relies on the following summation property for an $n$-th root of unity (where \(\omega^n = 1\)):
\(\sum_{k=0}^{n-1} \omega^{mk} = \begin{cases} n, & \text{if } m \text{ is a multiple of } n \\ 0, & \text{otherwise} \end{cases}\)\\

Given a polynomial function f(x), you can isolate the sum of every n-th coefficient (e.g., \(a_0, a_n, a_{2n}, \dots\)) by averaging f(x) evaluated at all n-th roots of unity:
\[ \sum_{k=0}^{\infty} a_{kn} = \frac{f(1) + f(\omega) + f(\omega^2) + \dots + f(\omega^{n-1})}{n} \]\\

%Applying It to Products
%In combinatorial problems (such as finding the number of subsets whose sum is divisible by a certain integer), your polynomial is often expressed as a product of factors, \(P(x) = \prod_{i=1}^{m} F_i(x)\). [1, 2]
%Applying the filter to this product involves substituting each root of unity \(\omega ^{j}\) into the product individually:\\

%Construct the Generating Function: Represent the event or choice as a polynomial, where powers indicate the "value" and coefficients indicate the number of combinations.
%Evaluate at Roots of Unity: Calculate \(P(\omega^j)\) for \(j = 0, 1, \dots, n-1\).
%Exploit the Product Rule: Because P(x) is a product, \(P(\omega^j)\) becomes the product of the individual evaluated factors:
%\(P(\omega^j) = F_1(\omega^j) \cdot F_2(\omega^j) \cdot \dots \cdot F_m(\omega^j)\) \\


%Average the Results: Sum these products and divide by n to isolate the terms matching your modulo condition. [1, 2, 3]\\

%Example: Divisibility of Subset Sums
%Suppose you want to find the number of subsets of the set \(\{1, 2, 3\}\) whose sum of elements is divisible by 2.

%Generating Function: Each element can be either included (x) or not included (1). The polynomial for the whole set is the product:
%P(x) = (1+x)(1+x^2)(1+x^3) [1]
%Apply the Filter (mod 2): Since we want the sum to be divisible by 2, we use the 2nd roots of unity, which are 1 and -1. [1]
%Evaluate and Add:
 %P(1) = (1+1)(1+1?)(1+1?) = 2 ? 2 ? 2 = 8
  %P(-1) = (1-1)(1+(-1)?)(1+(-1)?) = 0 ? 2 ? 0 = 0\\
  
 %Average:
 %\(\frac{P(1) + P(-1)}{2} = \frac{8 + 0}{2} = 4\)\\

%There are 4 valid subsets whose sum is divisible by 2: \(\{\emptyset\}, \{2\}, \{1,3\}, \{1,2,3\}\).\\

%Resources
%For a clear, visual explanation of how the roots of unity filter isolates terms in a product, check out this video:
